\chapter{Przegląd literatury i podstawy teoretyczne}\label{ch:przeglad}

\TemplateNote{Rozdział pokazuje cytowania, przypisy dolne, cytat blokowy oraz środowiska matematyczne
(definicja, twierdzenie, dowód, przykład, uwaga). Źródła znajdują się w pliku \file{bib/references.bib}.}

\section{Dotychczasowe prace}\label{sec:prace}

Skład tekstu przy pomocy systemu \LaTeX{} opisano szczegółowo w pracy~\cite{lamport1994}, a jego silnik
-- w klasycznym podręczniku~\cite{knuth1984}. Wiele wymagań dotyczących skrótów kryptograficznych
zebrano w normie~\cite{fips180-4}. Kilka źródeł cytuje się jednocześnie poleceniem
\latexcmd{cite\{a,b\}}, co daje zapis~\cite{lamport1994,knuth1984}.

Przypisy dolne tworzy polecenie \latexcmd{footnote}\footnote{Treść przypisu pojawia się na dole strony.}.
Dłuższe cytaty umieszcza się w środowisku \code{quote}:
\begin{quote}
    \enquote{Przykładowy cytat z literatury, ujęty w polskie cudzysłowy.}~\cite{knuth1984}
\end{quote}

\section{Podstawy teoretyczne}\label{sec:podstawy}

\begin{definition}\label{def:nieparzyste}
    Liczbę naturalną $k$ nazywamy \emph{nieparzystą}, jeżeli istnieje taka liczba $m \in \NN$, że $k = 2m + 1$.
\end{definition}

\begin{theorem}\label{thm:suma-nieparzystych}
    Dla każdej liczby naturalnej $n \geq 1$ suma $n$ początkowych liczb nieparzystych jest równa $n^2$:
    \begin{equation}\label{eq:suma-nieparzystych}
        \sum_{i=1}^{n} (2i - 1) = n^2 .
    \end{equation}
\end{theorem}

\begin{proof}
    Dowód indukcyjny. Dla $n = 1$ obie strony równości~\eqref{eq:suma-nieparzystych} są równe $1$. Załóżmy,
    że równość zachodzi dla pewnego $n$. Wtedy
    \[
        \sum_{i=1}^{n+1} (2i - 1) = n^2 + (2n + 1) = (n + 1)^2 ,
    \]
    co kończy dowód zgodnie z zasadą indukcji.
\end{proof}

\begin{example}
    Dla $n = 4$ otrzymujemy $1 + 3 + 5 + 7 = 16 = 4^2$, co zgadza się z twierdzeniem~\ref{thm:suma-nieparzystych}.
\end{example}

\begin{remark}
    Twierdzenia, definicje, przykłady i wnioski mają wspólny licznik numerowany w obrębie rozdziału.
\end{remark}

\section{Porównanie istniejących rozwiązań}\label{sec:porownanie}

Przegląd rozwiązań warto zakończyć krótkim zestawieniem (szczegóły tabel -- w rozdziale~\ref{ch:metodyka}):
\begin{enumerate}
    \item rozwiązanie pierwsze -- główne zalety i ograniczenia,
    \item rozwiązanie drugie:
    \begin{enumerate}
        \item podejście A,
        \item podejście B,
    \end{enumerate}
    \item rozwiązanie trzecie.
\end{enumerate}
